% Options for packages loaded elsewhere
\PassOptionsToPackage{unicode}{hyperref}
\PassOptionsToPackage{hyphens}{url}
%
\documentclass[
]{article}
\usepackage{amsmath,amssymb}
\usepackage{iftex}
\ifPDFTeX
  \usepackage[T1]{fontenc}
  \usepackage[utf8]{inputenc}
  \usepackage{textcomp} % provide euro and other symbols
\else % if luatex or xetex
  \usepackage{unicode-math} % this also loads fontspec
  \defaultfontfeatures{Scale=MatchLowercase}
  \defaultfontfeatures[\rmfamily]{Ligatures=TeX,Scale=1}
\fi
\usepackage{lmodern}
\ifPDFTeX\else
  % xetex/luatex font selection
\fi
% Use upquote if available, for straight quotes in verbatim environments
\IfFileExists{upquote.sty}{\usepackage{upquote}}{}
\IfFileExists{microtype.sty}{% use microtype if available
  \usepackage[]{microtype}
  \UseMicrotypeSet[protrusion]{basicmath} % disable protrusion for tt fonts
}{}
\makeatletter
\@ifundefined{KOMAClassName}{% if non-KOMA class
  \IfFileExists{parskip.sty}{%
    \usepackage{parskip}
  }{% else
    \setlength{\parindent}{0pt}
    \setlength{\parskip}{6pt plus 2pt minus 1pt}}
}{% if KOMA class
  \KOMAoptions{parskip=half}}
\makeatother
\usepackage{xcolor}
\usepackage[margin=1in]{geometry}
\usepackage{graphicx}
\makeatletter
\def\maxwidth{\ifdim\Gin@nat@width>\linewidth\linewidth\else\Gin@nat@width\fi}
\def\maxheight{\ifdim\Gin@nat@height>\textheight\textheight\else\Gin@nat@height\fi}
\makeatother
% Scale images if necessary, so that they will not overflow the page
% margins by default, and it is still possible to overwrite the defaults
% using explicit options in \includegraphics[width, height, ...]{}
\setkeys{Gin}{width=\maxwidth,height=\maxheight,keepaspectratio}
% Set default figure placement to htbp
\makeatletter
\def\fps@figure{htbp}
\makeatother
\setlength{\emergencystretch}{3em} % prevent overfull lines
\providecommand{\tightlist}{%
  \setlength{\itemsep}{0pt}\setlength{\parskip}{0pt}}
\setcounter{secnumdepth}{-\maxdimen} % remove section numbering
\ifLuaTeX
  \usepackage{selnolig}  % disable illegal ligatures
\fi
\usepackage{bookmark}
\IfFileExists{xurl.sty}{\usepackage{xurl}}{} % add URL line breaks if available
\urlstyle{same}
\hypersetup{
  pdftitle={Lab1},
  hidelinks,
  pdfcreator={LaTeX via pandoc}}

\title{Lab1}
\author{}
\date{\vspace{-2.5em}}

\begin{document}
\maketitle

\section{Empirical Exercise E2.1}\label{empirical-exercise-e2.1}

\subsection{Part A:}\label{part-a}

\begin{table}[h]
\centering
\begin{tabular}{cc}
\hline
Age & $P(\text{Age})$ \\
\hline
25 & 0.08489 \\
26 & 0.09223 \\
27 & 0.08547 \\
28 & 0.09339 \\
29 & 0.10350 \\
30 & 0.10473 \\
31 & 0.10393 \\
32 & 0.10807 \\
33 & 0.10880 \\
34 & 0.11498 \\
\hline
\end{tabular}
\caption{Marginal distribution of Age}
\end{table}

\subsection{Part B:}\label{part-b}

\begin{table}[h]
\centering
\begin{tabular}{cc}
\hline
Age & $E(AHE \mid \text{Age})$ \\
\hline
25 & 17.5908 \\
26 & 18.9669 \\
27 & 19.7049 \\
28 & 20.2358 \\
29 & 21.1714 \\
30 & 21.7849 \\
31 & 22.5951 \\
32 & 23.6920 \\
33 & 23.3487 \\
34 & 24.1081 \\
\hline
\end{tabular}
\caption{Conditional mean of hourly earnings given age}
\end{table}

\subsection{Part C:}\label{part-c}

\subsection{\texorpdfstring{\protect\includegraphics{Mean AHE for each Age vs. Age.png}}{}}\label{section}

\subsection{Part D:}\label{part-d}

Applying the Law of Iterated Expectations, \[
E(AHE)=\sum_{a} E(AHE \mid \text{Age}=a)\,P(\text{Age}=a).
\]

Using the conditional means and the marginal distribution of Age from
the spreadsheet, \[
E(AHE)=21.5106.
\]

The average hourly wage in the population is \textasciitilde\$21.51 from
the dataset.

\subsection{Part E:}\label{part-e}

The variance of average hourly earnings is given by \[
\mathrm{Var}(AHE)=E(AHE^2)-[E(AHE)]^2.
\]

\[
\mathrm{Var}(AHE)=164.8203.
\]

The standard deviation of hourly earnings is: \[
\sigma_{AHE}=\sqrt{164.8203}\approx 12.84.
\]

\subsection{Part F:}\label{part-f}

The covariance between \(AHE\) and Age is \[
\mathrm{Cov}(AHE,\text{Age})=E(AHE\cdot \text{Age})-E(AHE)\,E(\text{Age}).
\]

From the marginal distribution of Age, \[
E(\text{Age})=29.7663.
\]

\[
\mathrm{Cov}(AHE,\text{Age})=5.7172.
\]

Since the covariance is positive, higher ages tend to be associated with
higher hourly earnings.

\subsection{Part G:}\label{part-g}

The correlation between \(AHE\) and Age is \[
\mathrm{Corr}(AHE,\text{Age})
=\frac{\mathrm{Cov}(AHE,\text{Age})}
{\sqrt{\mathrm{Var}(AHE)}\sqrt{\mathrm{Var}(\text{Age})}}.
\]

From \[
\mathrm{Var}(AHE)=164.8203, 
\qquad 
\mathrm{Var}(\text{Age})=8.1447,
\] we get \[
\mathrm{Corr}(AHE,\text{Age})
=\frac{5.7172}{\sqrt{164.8203}\sqrt{8.1447}}
\approx 0.156.
\]

\subsection{Part H:}\label{part-h}

The covariance between age and average hourly earnings is positive,
which helps us recognize through data that earnings tend to increase
with age among workers aged 25--34. The correlation confirms a positive
linear relationship between age and earnings. This relationship is
visually supported by the upward-sloping pattern in the plot of
conditional mean earnings versus age.

\end{document}
